\section{Semantic twins}\label{sec:semantic}

VOUCH's secrets are random alphanumeric strings. A high-entropy string has no
semantically equivalent rephrasing, so a score that reads the secret under paraphrased
carrier sentences finds the same advantage as the literal score, and the squeezing
effect of \citet{li2026beliefs} cannot be observed; the original paper reports exactly
this and names natural-language secrets, exchangeable as propositions rather than as
strings, as the next step \citep{dang2026vouch}. This section defines such
canaries and shows that the exact null survives.

\subsection{The generator}

A semantic twin pair names a fresh fictitious entity $e$ and an attribute $a$ (for
example birthplace or profession) and states two incompatible values of it: $c^0$ says
that $e$ has value $v$, $c^1$ that $e$ has value $v'\neq v$. The pair is drawn by a
generator controlled by the verifier.

\begin{assumption}[Generator]\label{ass:gen}
The pair contents $C_1,\dots,C_m$ are independent draws from a generator distribution
$P$ on ordered pairs that is \emph{label-symmetric}: $(c^0,c^1)\sim P$ implies
$(c^1,c^0)\sim P$. The coins $b_i$ are independent $\Bern(\frac12)$ draws,
independent of the contents and of $\xi$, and $\xi$ is independent of the contents.
\end{assumption}

Both properties are under the verifier's control, because they describe its own
sampler rather than the data or the model. Our generator draws the two values as an
unordered pair from the attribute's value set and swaps them with probability one half,
which makes it label-symmetric by construction; the coins come from a separate random
stream. The in-twin is trained in one or more \emph{training forms}, for example the
literal statement ``$e$ was born in $v$'' and one question--answer form.

\subsection{Scores}

Each score compares the two values under the same prompt, so that both twins of a pair
are scored by the same function of the model:
\begin{itemize}
\item \emph{literal}: the mean token log-probability of the value after the literal
prefix (``$e$ was born in''), $s_{\mathrm{lit}}(M,c^k)=\overline{\log P_M}(v_k\mid
\text{prefix})$;
\item \emph{belief}: the same quantity averaged over a committed set $Q$ of paraphrased
questions (``Where was $e$ born?'', ``What is the hometown of $e$?''), none of which
appears in training;
\item \emph{belief\_max}: the maximum over $Q$ instead of the mean, the semantic
analogue of VOUCH's $s_{\mathrm{para}}$;
\item the \emph{calibrated} version of each, which subtracts the same score computed
under the frozen base model $M_0$: $s^{\mathrm{cal}}(M,c)=s(M,c)-s(M_0,c)$.
\end{itemize}
Calibration removes the base model's prior preference between the two values (a base
model may find ``born in Paris'' more likely than ``born in Porto'' for any name). It
keeps the score in the declared class because $M_0$ is fixed before the coins are drawn.
In the experiments we refer to the literal score as \emph{literal} and to the calibrated
belief scores as \emph{paraphrase}.

\subsection{The exact null}

\begin{proposition}[Exact null for semantic twins]\label{prop:semantic}
Under Assumption~\ref{ass:gen}, if the deployed model is conditionally independent of
the coins given the contents, then for every score $s$ of the form above and every step
$t$ of a reveal order chosen independently of the signs,
$Z_{\pi(t)}^{(s)}\mid\F_{t-1}\sim\Bern(\frac12)$ exactly.
\end{proposition}

\begin{proof}
The proof of VOUCH Theorem~1 uses only that, given the unordered pair opened at step
$t$, the two numbers $s(M,c^0)$ and $s(M,c^1)$ are determined, and that the coin
deciding which of them belongs to the in-twin is fair given everything opened before
and independent of the model given the contents. Neither the semantics of the records
nor the form of $s$ enters, provided $s(M,\cdot)$ is a function of the model and the
record alone. Each score above is such a function: both twins share the prefix or the
question set and differ only in the value evaluated, and the calibrated score subtracts
a function of $M_0$ and the record, where $M_0$ does not depend on the coins.
\end{proof}

Two remarks delimit what the proposition gives. First, \emph{semantics costs power, not
validity}. A base-model prior between the two values enters both arms symmetrically and
leaves the null exact, but it adds variance to the score differences; calibration
removes most of it. Second, the null does not need label symmetry: the coin, not the order of the pair,
decides which twin is trained. Label symmetry matters for what the certificate says
about records outside the cohort. It makes the two components $c^0$ and $c^1$ of a pair
identically distributed, so every canary record is a draw from a single record
distribution, whichever position it occupies and whichever twin the coin trains, and it is this distribution that Section~\ref{sec:organic} compares with
the records users actually delete.
